\documentclass[12pt,a4paper]{book}
\newcommand{\niveau}{6\textsuperscript{e}}
\input{../../commun/preambule_commun}

\title{Cours de Mathématiques — Classe de 6\textsuperscript{e}\\[0.4em]\large Année scolaire 2025–2026}
\author{Abdoullatuf Maoulida}
\date{\today}


% Commande pour importer une séquence en forçant le numéro de chapitre
% Usage : \InputSeq{Numéro}{Chemin/Fichier}
\newcommand{\InputSeq}[2]{%
	\setcounter{chapter}{\numexpr#1-1\relax}%
	\input{#2}%
}

\begin{document}
% VERSION ÉLÈVE (par défaut - avec espaces à trous)
\maketitle
\tableofcontents
\cleardoublepage

%\InputSeq{1}{chapitres/seq_01_nombres_entiers} % Les nombres entiers
%\InputSeq{2}{chapitres/seq_02_premiers_elts_geometrie} % Points et droites
%\InputSeq{3}{chapitres/seq_03_fractions_decimales_et_nombres_decimaux} % Fractions décimales et nombres décimaux
%\InputSeq{4}{chapitres/seq_04_distances_et_cercles}
%\InputSeq{5}{chapitres/seq_05_cercles_et_triangles_partie_2}
\InputSeq{6}{chapitres/seq_06_proportionnalite_partie_1}
% \InputSeq{16}{chapitres/seq_16_proportionnalite_partie_2}

\cleardoublepage
\appendix
\chapter{Progression annuelle (récapitulatif)}
Cette progression correspond à la répartition établie pour l'année 2025–2026.

\begin{center}
\renewcommand{\arraystretch}{1.2}
\begin{tabularx}{\textwidth}{|>{\raggedright\arraybackslash}p{0.32\textwidth}|>{\raggedright\arraybackslash}X|}
\hline
\textbf{Période} & \textbf{Séquences}\\ \hline
Période 1 (6 semaines) & S01 -- Les nombres entiers, S02 -- Points et droites, S03 -- Fractions décimales et nombres décimaux\\ \hline
Période 2 (7 semaines) & S04 -- Distance, cercle et triangles, S05 -- Notion de proportionnalité, S06 -- Notion de probabilités, S07 -- Angles et rapporteur\\ \hline
Période 3 (6 semaines) & S08 -- Opérations avec les nombres décimaux, S09 -- La médiatrice d'un segment, S10 -- La division, S11 -- Symétrie axiale\\ \hline
Période 4 (7 semaines) & S12 -- Fraction partage et comparaison de fractions, S13 -- Unités de longueur, de masse et de contenance, S14 -- Calculer avec les angles, S15 -- Nombres en écriture fractionnaire\\ \hline
Période 5 (6 semaines) & S16 -- Proportionnalité et pourcentages, S17 -- Déterminer des probabilités et des issues, S18 -- Aires et périmètres, S19 -- Heures et durées, S20 -- Solides et volumes\\ \hline
\end{tabularx}
\end{center}


\end{document}
